% Generated draft: requires mathematical and source review.
\documentclass[11pt]{article}
\usepackage[margin=25mm]{geometry}
\usepackage{fontspec}
\setmainfont{Latin Modern Roman}
\usepackage{amsmath,amssymb,mathtools,unicode-math}
\setmathfont{Latin Modern Math}
\usepackage{enumitem,needspace,xurl,hyperref}
\setlength{\parindent}{0pt}
\setlength{\parskip}{5pt}
\setlength{\emergencystretch}{3em}
\newcommand{\N}{\mathbb{N}}
\newcommand{\Z}{\mathbb{Z}}
\newcommand{\Q}{\mathbb{Q}}
\newcommand{\R}{\mathbb{R}}
\newcommand{\C}{\mathbb{C}}
\newcommand{\F}{\mathbb{F}}
\newcommand{\A}{\mathbb{A}}
\newcommand{\PP}{\mathbb{P}}
\newcommand{\E}{\mathbb{E}}
\newcommand{\eps}{\varepsilon}
\newcommand{\dd}{\,\mathrm{d}}
\begin{document}
\section*{UnsolvedMath ID: 700309}
\textbf{Source problem code:} catalogue entry
\subsection{Definitions and mathematical statement}
A group \(G\) is soluble if its derived series, defined by \(G^{(0)}=G\) and \(G^{(n+1)}=[G^{(n)},G^{(n)}]\), satisfies \(G^{(r)}=\{1\}\) for some integer \(r\geq 0\). Here \([H,H]\) is the subgroup generated by all commutators \(x^{-1}y^{-1}xy\) with \(x,y\in H\), and \(G'=[G,G]\). Under the usual two-sided convention, a linearly ordered group is a group equipped with a total order \(\leq\) satisfying \(x\leq y\Longrightarrow axb\leq ayb\) for all \(a,b,x,y\in G\). A group is periodic if every element has finite order; it is non-periodic if it contains an element of infinite order.

The supplied assertion is that, for every soluble linearly ordered group \(G\), the quotient \(G/G'\) is non-periodic. Equivalently, the requested conclusion is
\[
\exists g\in G\quad\forall n\in\Z_{\geq 1},\qquad g^n\notin G'.
\]
The source does not explicitly exclude the trivial group. With the definitions above, the trivial group satisfies the hypotheses but has periodic abelianization, so a nontriviality qualification is necessary for the intended assertion. Whether \(G\neq\{1\}\) is an implicit source convention requires review. The supplied record also does not explicitly specify the sidedness of order invariance; the definition above states the usual two-sided interpretation.
\subsection{Short English statement}
V. M. Kopytov asks whether the abelianization of every soluble linearly ordered group contains an element of infinite order. The intended exclusion of the trivial group and the convention for order invariance require confirmation from the source.
\subsection{Sources}
\begin{itemize}
\item Local UnsolvedMath registry: permanent ID \(700309\), exact source code \texttt{catalogue entry}, title \emph{Kourovka Notebook Problem 5.25}. The supplied registry labels the problem open; this status is not independently verified.
\item V. M. Kopytov, problem 5.25, \emph{The Kourovka Notebook (Unsolved Problems in Group Theory)}, 5th issue (1976), as attributed in the supplied record.
\item The supplied background identifies the 21st issue edition (2026), arXiv:1401.0300v48. The registry describes the Issues 1--21 collection as edited by E. I. Khukhro and V. D. Mazurov. No external source was consulted.
\end{itemize}
\end{document}
